\section{Discussion}\label{sec:discussion}

Within the two-dimensional algorithm of Theorem~\ref{thm:main}, every scalar and polynomial multiplication across the computation (root pointwise squaring, overlap-add Kronecker spreading, pre- and post-twists, two-term and Nyquist column deconvolution, and leaf radix-2 FFT butterflies) executes in $O\bigl(n(\log\log n)^2\bigr)$ steps on a multitape Turing machine, while shared per-level table generation takes $o(n)$ steps. The only $\Theta(n \log n)$ term in the upper bound is $\mathcal{C}_{\mathrm{syn+route}}$ in~\eqref{eq:syn-route-total}: the multiplication-free polynomial transform $G_{t_{1,\ell}}$ over $\C[y]/(y^{r_\ell} + 1)$ together with the per-level matrix transposes and odd-stride tape routing. Known $\Omega(n \log n)$ lower bounds for integer multiplication are conditional, resting on the network-coding conjecture (for Boolean circuits)~\cite{AFKL19} or on the conjectured $\Omega(n^2 \log n)$ multitape Turing machine complexity of $n \times n$ binary matrix transposition~\cite{HvdH25}. Whether the $\Theta(n \log n)$ polynomial-transform and tape-routing term in a fixed-dimension algorithm can be reduced unconditionally remains open, although~\cite{OAI26} gives, and~\cite{CLW26} claims, $o(n \log n)$ bounds by letting the dimension grow with $n$, and no lower bound is claimed in this paper.
